\section{Introduction}\label{sec:intro}

Let $Q$ be an irreducible generator on $[d]=\{1,\dots,d\}$, and let $J_V$ be the
Donsker and Varadhan functional of its occupation measure~\cite{dv1975}. We prove
that $J_V$ determines exactly the principal minors of $-Q$
(Theorem~\ref{thm:A}). The generators with these minors form a finite set. It
consists of $Q$ and its time reversal $Q^*$ when $-Q$ has no rank-one cut, and it is
an orbit of cut-transposes in general. Kim and Landim~\cite[Section~5]{kl2024} show
that $J_V$ determines the minors of orders 1 and 2, and they observe that the
uniform 3-cycle and its reversal share $J_V$. Theorem~\ref{thm:A} says exactly
which generators share it.

The functional enters through the occupation counts of a chain that rarely
switches. Fix a start law and run the chain for $N$ steps with transition matrix
$I+(T/N)Q$. With $T=\tau\log N$ the chain makes about $\tau\log N$ switches. Its
composition is the vector of occupation counts, and the support of a composition is
a face of the simplex. In~\cite{paperC1} we describe the law of the composition on
every face at this scale. On a face $F$ the most likely composition has probability
$N^{-\Phi_\tau(F)+o(1)}$, with the score $\Phi_\tau(F)=|F|-1+\tau\lambda_F$ and
$\lambda_F$ the principal exit rate of $F$. One power of $N$ per dimension competes
with the cost of staying in the face, and the mode moves to larger faces as $\tau$
grows. Inside a face the logarithm of the law is $T$ times $J_V$.

We ask which coarse functionals of the law of the counts determine $Q$, and what
each of them loses. We compare 3, from fine to coarse. The first is the law itself,
at every $N$ and $T$. The second is its logarithmic scale on single compositions,
the functional $J_V$. The third is a finite label, the support of the most likely
composition when the chain is killed at known rates. The third determines exactly
the same minors as the second (Theorem~\ref{thm:B}). We then read a reversible
generator from finitely many boundaries of a design fixed in advance
(Theorem~\ref{thm:C}), and the Freidlin and Wentzell cycle tree of an Arrhenius
family from the same labels (Theorem~\ref{thm:D}).

\subsection{The Donsker and Varadhan functional}\label{sec:intro-dv}

Write $Q^*=D_\pi^{-1}Q^\top D_\pi$ for the time reversal of an irreducible
generator with stationary law $\pi$. Call a partition $[d]=X\cup Y$ with
$|X|,|Y|\ge2$ a cut of $Q$ if both blocks $Q[X,Y]$ and $Q[Y,X]$ have rank at most
$1$. Across a cut a cut-transpose replaces the 2 rank-one blocks by the products
of their factors taken crosswise and transposes the block on $Y$ (Section~\ref{sec:fiber}).

\begin{introtheorem}[what the Donsker and Varadhan functional determines]\label{thm:A}
Let $Q$ and $Q'$ be irreducible generators on $[d]$.
\begin{enumerate}[(i)]
\item $J_V=J'_V$ on the simplex if and only if $-Q$ and $-Q'$ have the same
principal minors.
\item The irreducible generators with the principal minors of $-Q$ form a finite
set. It contains $Q$ and $Q^*$. It is $\{Q,Q^*\}$ when $Q$ has no cut, in
particular when $d\le3$, and $\{Q\}$ when $Q$ is reversible. In general it is the
set of generator points of the cut-transpose orbit of $-Q$, and it can contain
generators other than $Q$ and $Q^*$.
\item From the stationary start the law of the composition is the same under $Q$
and under $Q^*$, for every $N$ and every $T$ with $Tq_{\max}\le N$.
\end{enumerate}
\end{introtheorem}

Part (i) is Theorem~\ref{thm:dv}, part (ii) is Theorem~\ref{thm:nonreversible}
with the example after it, and part (iii) is Proposition~\ref{prop:stationary}.
Part (i) has 3 inputs. The first is the variational formula
$\lambda_V(v)=\min_\alpha(J_V(\alpha)-v\cdot\alpha)$ for
$\lambda_V(v)=-\varrho(Q+\operatorname{diag}v)$. The second is the irreducibility
of the multiaffine polynomial $\det(-Q_F-\operatorname{diag}w)$ for every
irreducible face, which is Lemma~5.1 of Al Ahmadieh~\cite{alahmadieh2023}. The third
is the Perron-branch lemma (Lemma~\ref{lem:perron-branch}). It says that the
Perron root on an open set of fields determines that polynomial, hence every
principal minor. Kim and Landim use the minors of orders 1 and 2, with the
stationary law, to fix a reversible generator. By their Corollary~5.8 the joint
functional of the empirical measure and the empirical flow determines the
generator, so adding the flow removes the whole fiber. Part (iii) shows that the
fiber is not an artefact of the logarithmic scale. From the stationary start even
the exact law of the counts cannot separate $Q$ from $Q^*$. For a start outside a
finite union of proper subspaces the second-order amplitudes of the local law do
separate every member of the fiber (Proposition~\ref{prop:amplitudes}).

\subsection{Modal supports}\label{sec:intro-modal}

Weight each composition by $e^{Tv\cdot n/N}$ for a field $v\le0$. This is the
chance that the chain survives when it is killed at state $i$ at the known rate
$-v_i$. We record only which face carries the most likely surviving composition,
as a function of $(\tau,v)$, and call this the modal-support map. The data are the
boundaries between the open regions where one face wins strictly
(Definition~\ref{def:controlled}).

\begin{introtheorem}[what the modal supports determine]\label{thm:B}
Let $Q$ and $Q'$ be irreducible generators and let the start have full support.
They have the same modal-support map, outside a closed nowhere dense set of
$(\tau,v)$, if and only if $-Q$ and $-Q'$ have the same principal minors. If the
minors differ, there is a nonempty open set of $(\tau,v)$ on which both maps have
a strict winner and the 2 winners differ.
\end{introtheorem}

This is Theorem~\ref{thm:modal}, with Proposition~\ref{prop:identifiability} for
the last sentence. The proof is a descent through the faces. For every irreducible
face $F$ and a generic field on $F$, the exposure lemma (Lemma~\ref{lem:exposure})
finds a proper subface $G$ that shares a boundary of strict regions with $F$. Along
that boundary $\lambda_F=\lambda_G-(|F|-|G|)/\tau$. Induction on $|F|$ then traces
$\lambda_F$ on an open set of fields, and the Perron-branch lemma extends it to all
fields.

Theorem~\ref{thm:B} is a statement about retention, not about sampling. The
survival probability under calibrated killing is a finer observation. It equals
$e^{-T\lambda_V(v)}$ up to a bounded factor, so it gives every principal minor at
once, with no descent (Remark~\ref{rem:survival}). The modal support is not easier
to sample. The most likely composition carries probability $N^{-\Phi+o(1)}$, so
estimating which face carries it needs far more independent runs than estimating
the survival probability. We know of no experiment in which the modal support is
observed and the survival probability is not. But the forward theory
of~\cite{paperC1} outputs the support of the most likely composition, and so does
any exact optimizer that reports only a maximizer of the weighted law.
Theorem~\ref{thm:B} says that this output loses nothing that $J_V$ keeps.

\emph{What is borrowed and what is new.} The principal minors are not observed.
Theorem~\ref{thm:dv} reads all of them off $J_V$, and Theorem~\ref{thm:modal} reads
all of them off the modal-support map, through the exposure lemma and the
Perron-branch lemma. This is our main contribution. The fiber of the
principal-minor map of an irreducible matrix is known, up to diagonal similarity and
transposition, from Loewy~\cite{loewy1986}, Hartfiel and
Loewy~\cite{hartfielloewy1984}, Al Ahmadieh~\cite{alahmadieh2023} and Chatterjee,
Ghosh, Gurjar and Raj~\cite{cggr2024}, who describe it by sequences of
cut-transposes (Section~\ref{sec:fiber}). Our second contribution, and the shortest
step, translates this fiber to generators. A Perron
and Frobenius normalization picks exactly one generator in each class of diagonal
similarity. So the generator fiber is finite, every member is irreducible with the
spectrum of $Q$, and a cut-transpose acts as $Q$ on one side of the cut and as $Q^*$
on the other (Theorem~\ref{thm:nonreversible} and
Proposition~\ref{prop:amplitudes}). The third is the stationary obstruction of
part (iii) of Theorem~\ref{thm:A}.

\subsection{A reversible design and finite horizons}\label{sec:intro-reversible}

For a reversible chain on a known connected graph with $m$ edges the fields can be
fixed before anything is observed.

\begin{introtheorem}[a reversible design]\label{thm:C}
Let $Q$ be reversible and irreducible on a known connected graph with $m$ edges,
and suppose that known constants $R$ and $b$ bound the norm of the symmetrized
generator above and every edge product $\sqrt{q_{ij}q_{ji}}$ below.
\begin{enumerate}[(i)]
\item A design whose fields and time scales depend only on $R$, $b$ and the
labels reads $d-1$ singleton switches and $m$ pair-entry times off the
modal-support map, and these $m+d-1$ numbers determine $Q$.
\item No family of fewer than $m+d-1$ scalar $C^1$ observations, with finitely many
locally constant labels, identifies $Q$ on any open set of such generators.
\item The reconstruction is Lipschitz in the observed coordinates, with constants
that depend only on $d$ and $b/R$.
\item At the finite horizon $N$, with $L=\log N$, the observed pair-entry times are
$t_{ij}(1-\log L/(2L))+c_{ij}/L+O(\log L/L^2)$, uniformly over the generators of
the class and the starts with $\min_i\mu_i\ge\mu_*$, and the start bound cannot be
dropped.
\end{enumerate}
\end{introtheorem}

Parts (i) to (iii) are Theorems~\ref{thm:reversible-reconstruction},
\ref{thm:minimality} and~\ref{thm:stability}, and part (iv) is
Theorem~\ref{thm:finite-horizon} with Remark~\ref{rem:finite-start}. The coordinates
are those of Wolfer and Watanabe~\cite{wolferwatanabe2021}, and Kim and
Landim~\cite[Lemma~5.1 and Corollary~5.3]{kl2024} recover the same quantities from
$J_V$. Theorem~\ref{thm:C} adds a design fixed in advance, its minimality, its
stability with constants from $R$ and $b$ alone, and its finite-horizon behaviour.

\subsection{Low temperature and the cycle tree}\label{sec:intro-lowtemp}

Let the rates be $q_{ij}=e^{-V_{ij}/\eta}$ with fixed costs $V_{ij}\ge0$ and a
temperature $\eta\downarrow0$, and let the time scale be $\tau=e^{h/\eta}$ for a
height $h$. Then dimension and exit rate compete on the Freidlin and Wentzell cycles
of the chain. The depth $D(F)$ of a face is the exponential order of its exit rate,
and it is the largest depth of a cycle inside $F$ (Theorem~\ref{thm:depth}). Away
from finitely many heights the winner is a cycle of least size among those deeper
than $h$ (Theorem~\ref{thm:noncritical}). At a critical height explicit prefactors
decide, and a noncycle can win on an open window of the rescaled time scale
(Theorem~\ref{thm:critical-window} and Example~\ref{ex:four-state}).

So modal winners need not be cycles. Do the modal supports determine the cycles at
all? The 2 obvious experiments do not. The unperturbed modal path shows the depth of
each face it selects and nothing about the cycles that are never the smallest
eligible face. Killing the complement of a face at a large rate does not show its
depth either. The whole space then has an exit rate of the same order, and the
winner leaves the face only at a larger height. The cure is to reward the face
instead. For a proper face $G$, a constant $c>0$ and a real $\gamma$ we put the field
$v^{G,\gamma}=c\,e^{-\gamma/\eta}\ones_G$. The threshold $\vartheta(G)$ is the
supremum of the $\gamma$ for which $G$ wins under $v^{G,\gamma}$ at some height above
$\gamma$.

\begin{introtheorem}[the cycle tree from calibrated rewards]\label{thm:D}
Let $(Q_\eta)$ be an Arrhenius family with all rates positive, and let the start have
full support.
\begin{enumerate}[(i)]
\item $\vartheta(G)\le D(G)$ for every proper face $G$, with equality when $G$ is a
cycle.
\item $D(F)=\max_{G\subseteq F}\vartheta(G)$ for every proper face $F$.
\item A nonempty proper set $C$ is a cycle if and only if $D(C')<D(C)$ for every
nonempty proper subset $C'$ of $C$.
\end{enumerate}
So the modal-support map, restricted to the fields $v^{G,\gamma}$, determines the
depth of every face, the cycles, their depths and the cycle tree. No step uses
reversibility, and no nondegeneracy of the depths is assumed.
\end{introtheorem}

Parts (i) and (ii) are Theorem~\ref{thm:attraction} and part (iii) is
Theorem~\ref{thm:cycles-from-depths}. The design needs 3 calibrated quantities,
the time unit, the amplitude of the reward and the labels of the states, and no rate
of $Q$. The reward keeps a face whose rate is known, namely the whole space with
$\lambda_V(v)\le0$, and the killing design destroys exactly this. Criterion (iii)
holds with ties between the depths of different faces, because Catoni's cycles are
strict and depths increase strictly along nested cycles (Lemma~\ref{lem:nested}). Its
weak form, with $D(C')\le D(C)$, fails already on 3 states. Finally the endpoints
of the critical windows determine the prefactors of the faces on the lower envelope,
and only bound the others (Theorem~\ref{thm:window-prefactors}).

\subsection{Related work}\label{sec:related}

The rate inside a face is the Donsker and Varadhan functional of the killed
chain~\cite{dv1975}, and the exact law of the counts goes back to
Good~\cite{good1961}. Bertini, Gabrielli and Landim~\cite{bgl2024} and
Landim~\cite{landim2023} expand the level-two rate functions of finite chains in a
parameter, and Kim and Landim~\cite{kl2024} treat the joint rate of the empirical
measure and the flow~\cite{bfg2014,bfg2015}. Their Section~5 is the starting point of
Theorem~\ref{thm:A}. The matrix theory behind the fiber is credited above and in
Section~\ref{sec:fiber}, together with Engel and Schneider~\cite{engelschneider1980}.
Cohen~\cite{cohen1981} and Friedland~\cite{friedland1981} study the convexity of
$v\mapsto\varrho(A+\operatorname{diag}v)$, and Delmas, Dronnier and Zitt~\cite{ddz2024}
prove a multiplicative analogue of the Perron-branch lemma. The cycle framework is that of
Freidlin and Wentzell~\cite{freidlinwentzell2012}, Catoni~\cite{catoni1999}, Olivieri and
Vares~\cite{olivierivares2005}, Trouv\'e~\cite{trouve1996} and Cirillo, Nardi and
Sohier~\cite{cirillonardisohier2015}. The eigenvalue asymptotics behind the prefactors go
back to Wentzell~\cite{wentzell1972}. New at low temperature are the score of dimension
against exit rate, the competition of envelopes at a critical depth with its noncycle window,
and the reading of the cycle tree from calibrated rewards.

\subsection{Plan of the paper}\label{sec:plan}

Section~\ref{sec:prelim} fixes the notation and states the results of~\cite{paperC1}
that we use. They are the first-order results, the face maxima with a field, their
uniform version on a compact class and the constant of the local law.
Section~\ref{sec:minors} proves Theorem~\ref{thm:A}. Section~\ref{sec:modal} sets up
the modal-support map and proves Theorem~\ref{thm:B}. Section~\ref{sec:reversible}
treats the reversible design and the finite horizon, and Section~\ref{sec:lowtemp} the
low temperature. Section~\ref{sec:verify} lists the computations that check the paper
and the open questions.
